\documentclass[a0paper, portrait, fontscale=0.38]{tikzposter}

%%%%%%%%%%%%%%%%%%%%%%%%%%%%% Paquetes %%%%%%%%%%%%%%%%%%%%%%%%%%%%%%%%%%%%%%%%
\usepackage[spanish]{babel}
\usepackage{amsmath}
\usepackage{amssymb}
\usepackage{booktabs}
\usepackage{graphicx}
\usepackage{pgfplots}
\usepackage{multicol}
\usepackage{apacite}
\pgfplotsset{compat=1.9}
%%%%%%%%%%%%%%%%%%%%%%%%%%%%%%%%%%%%%%%%%%%%%%%%%%%%%%%%%%%%%%%%%%%%%%%%%%%%%%%

%%%%%%%%%%%%%%%%%%%%%%%%%%%%% Paleta de colores %%%%%%%%%%%%%%%%%%%%%%%%%%%%%%%
\definecolor{fondobase}{RGB}{10, 15, 35}
\definecolor{acento}{RGB}{0, 180, 216}
\definecolor{acentoclaro}{RGB}{144, 224, 239}
\definecolor{bloque}{RGB}{20, 28, 58}
\definecolor{textoclaro}{RGB}{220, 230, 240}
\definecolor{blancopuro}{RGB}{255, 255, 255}
%%%%%%%%%%%%%%%%%%%%%%%%%%%%%%%%%%%%%%%%%%%%%%%%%%%%%%%%%%%%%%%%%%%%%%%%%%%%%%%

%%%%%%%%%%%%%%%%%%%%%%%%%%%%% Estilo tikzposter %%%%%%%%%%%%%%%%%%%%%%%%%%%%%%%
\definecolorstyle{UISstyle}{}{
    \colorlet{colorOne}{fondobase}
    \colorlet{colorTwo}{acento}
    \colorlet{colorThree}{bloque}
}

\definebackgroundstyle{UISbg}{
    \fill[fondobase]
        (current page.south west) rectangle (current page.north east);
    \foreach \i in {1,...,12}{
        \draw[acento!10, line width=0.6pt]
            ([xshift=\i*8cm]current page.south west) --
            ([xshift=\i*8cm]current page.north west);
    }
    \draw[acento, line width=4pt]
        (current page.south west) -- (current page.south east);
    \draw[acento, line width=4pt]
        (current page.north west) -- (current page.north east);
}

\definetitlestyle{UIStitle}{
    width=\textwidth, linewidth=4pt, titletotopverticalspace=2cm,
    titletoblockverticalspace=2cm
}{}

\defineblockstyle{UISblock}{
    titlewidthscale=1, bodywidthscale=1,
    titleleft, titleoffsetx=0pt, titleoffsety=0pt,
    bodyoffsetx=0pt, bodyoffsety=0pt,
    bodyverticalshift=0pt, roundedcorners=6,
    linewidth=1.5pt, titleinnersep=8mm, bodyinnersep=8mm
}{
    \fill[bloque, rounded corners=6pt] (blockbody.south west)
        rectangle (blocktitle.north east);
    \draw[acento, line width=1.5pt, rounded corners=6pt]
        (blockbody.south west) rectangle (blocktitle.north east);
    \draw[acento, line width=2pt]
        (blocktitle.south west) -- (blocktitle.south east);
}

\usetheme{Default}
\usecolorstyle{UISstyle}
\usebackgroundstyle{UISbg}
\usetitlestyle{UIStitle}
\useblockstyle{UISblock}

\colorlet{titlebgcolor}{fondobase}
\colorlet{titlefgcolor}{blancopuro}
\colorlet{blocktitlebgcolor}{bloque}
\colorlet{blocktitlefgcolor}{acento}
\colorlet{blockbodybgcolor}{bloque}
\colorlet{blockbodyfgcolor}{textoclaro}
\colorlet{innerblocktitlebgcolor}{fondobase}
\colorlet{innerblocktitlefgcolor}{acentoclaro}
\colorlet{innerblocktitlebgcolor}{fondobase}
\colorlet{notebgcolor}{fondobase}
\colorlet{notefrcolor}{acento}
%%%%%%%%%%%%%%%%%%%%%%%%%%%%%%%%%%%%%%%%%%%%%%%%%%%%%%%%%%%%%%%%%%%%%%%%%%%%%%%

%%%%%%%%%%%%%%%%%%%%%%%%%%%%% Título %%%%%%%%%%%%%%%%%%%%%%%%%%%%%%%%%%%%%%%%%%
\title{\parbox{0.85\textwidth}{\centering
    Control MPPT en Convertidor Boost para una Carga de 8540 W
}}
\author{Daniel Aranda Contreras \quad Dairo Lobo Moreno \quad Yulieth Portilla Jaimes}
\institute{Escuela E3T · Universidad Industrial de Santander · Bucaramanga, Colombia}
\titlegraphic{} % puedes poner \includegraphics[height=6cm]{logo.png}
%%%%%%%%%%%%%%%%%%%%%%%%%%%%%%%%%%%%%%%%%%%%%%%%%%%%%%%%%%%%%%%%%%%%%%%%%%%%%%%

\begin{document}
\maketitle

%%%%%%%%%%%%%%%%%%%%%%%%%%%%% Fila 1 %%%%%%%%%%%%%%%%%%%%%%%%%%%%%%%%%%%%%%%%%%
\begin{columns}
\column{0.33}

\block{1. Introducción}{
    Los sistemas fotovoltaicos requieren estrategias de control que maximicen
    la potencia extraída independientemente de las condiciones ambientales.
    El algoritmo \textbf{Perturba y Observa (P\&O)} es ampliamente usado por
    su simplicidad y efectividad \cite{mohan2003}.

    \vspace{0.5cm}
    \textbf{Objetivos:}
    \begin{itemize}
        \item Diseñar el convertidor Boost para la carga especificada.
        \item Implementar control MPPT con algoritmo P\&O.
        \item Validar mediante simulación en Simulink.
    \end{itemize}
}

\block{2. Sistema Fotovoltaico}{
    \begin{tikzpicture}
        \node[draw=acento, fill=bloque, rounded corners=4pt,
              text=textoclaro, minimum width=3cm, minimum height=1.2cm,
              font=\small] (panel) at (0,0) {Panel FV};
        \node[draw=acento, fill=bloque, rounded corners=4pt,
              text=textoclaro, minimum width=3cm, minimum height=1.2cm,
              font=\small] (boost) at (5,0) {Boost};
        \node[draw=acento, fill=bloque, rounded corners=4pt,
              text=textoclaro, minimum width=3cm, minimum height=1.2cm,
              font=\small] (carga) at (10,0) {Carga};
        \node[draw=acentoclaro, fill=fondobase, rounded corners=4pt,
              text=textoclaro, minimum width=3cm, minimum height=1.2cm,
              font=\small] (mppt) at (5,-2.5) {MPPT P\&O};
        \draw[->, acento, thick] (panel) -- (boost);
        \draw[->, acento, thick] (boost) -- (carga);
        \draw[->, acentoclaro, thick, dashed] (mppt) -- (boost);
        \draw[->, acentoclaro, thick, dashed] (boost.south) |- (mppt.east);
    \end{tikzpicture}
}

\column{0.34}

\block{3. Modelo Matemático}{
    Tensión de salida del convertidor Boost en MCC:
    \begin{equation}
        V_o = \frac{V_{in}}{1-D}, \qquad D \in [0,1)
    \end{equation}

    Rizado de corriente y tensión:
    \begin{equation}
        \Delta i_L = \frac{V_{in} \cdot D}{f_s \cdot L}, \qquad
        \Delta V_o = \frac{I_o \cdot D}{f_s \cdot C}
    \end{equation}

    \vspace{0.3cm}
    \begin{center}
    \begin{tabular}{lcc}
        \toprule
        \textbf{Parámetro} & \textbf{Símbolo} & \textbf{Valor} \\
        \midrule
        Tensión entrada   & $V_{in}$  & 200 V   \\
        Tensión salida    & $V_o$     & 400 V   \\
        Potencia          & $P_o$     & 8540 W  \\
        Inductancia       & $L$       & 2.3 mH  \\
        Capacitancia      & $C$       & 470 µF  \\
        Frecuencia        & $f_s$     & 20 kHz  \\
        \bottomrule
    \end{tabular}
    \end{center}
}

\block{4. Respuesta en Simulación}{
    \begin{tikzpicture}
        \begin{axis}[
            width=0.9\linewidth, height=7cm,
            xlabel={\textcolor{textoclaro}{Tiempo [ms]}},
            ylabel={\textcolor{textoclaro}{Tensión [V]}},
            grid=major,
            grid style={dotted, acento!30},
            axis background/.style={fill=fondobase},
            tick style={color=textoclaro},
            ticklabel style={color=textoclaro, font=\small},
            label style={color=textoclaro},
            legend style={fill=bloque, draw=acento, text=textoclaro, font=\small},
            legend pos=south east,
        ]
            \addplot[acento, thick, domain=0:10, samples=200]
                {400 - 380*exp(-x/1.5)};
            \addplot[acentoclaro, thick, dashed, domain=0:10, samples=200]
                {200 + 5*sin(deg(2*pi*x))};
            \legend{$V_o$, $V_{in}$}
        \end{axis}
    \end{tikzpicture}
}

\column{0.33}

\block{5. Algoritmo P\&O}{
    \begin{tikzpicture}[
        every node/.style={font=\small},
        box/.style={draw=acento, fill=bloque, rounded corners=4pt,
                    text=textoclaro, minimum width=4.5cm,
                    minimum height=0.9cm, align=center},
        dec/.style={draw=acentoclaro, fill=fondobase, diamond,
                    text=textoclaro, minimum width=1.2cm,
                    minimum height=1.2cm, align=center, font=\scriptsize}
    ]
        \node[box] (med)  at (0,0)    {Medir $V$, $I$};
        \node[box] (pot)  at (0,-2)   {Calcular $P = VI$};
        \node[dec] (dp)   at (0,-4)   {$\Delta P > 0$};
        \node[dec] (dv)   at (-3,-6)  {$\Delta V > 0$};
        \node[dec] (dv2)  at (3,-6)   {$\Delta V > 0$};
        \node[box] (inc)  at (-4.5,-8.5) {$D = D - \Delta D$};
        \node[box] (dec1) at (-1.5,-8.5) {$D = D + \Delta D$};
        \node[box] (dec2) at (1.5,-8.5)  {$D = D - \Delta D$};
        \node[box] (inc2) at (4.5,-8.5)  {$D = D + \Delta D$};

        \draw[->, acento, thick] (med)  -- (pot);
        \draw[->, acento, thick] (pot)  -- (dp);
        \draw[->, acento, thick] (dp)   -- node[left,font=\scriptsize,text=textoclaro]{Sí} (dv);
        \draw[->, acento, thick] (dp)   -- node[right,font=\scriptsize,text=textoclaro]{No} (dv2);
        \draw[->, acentoclaro]   (dv)   -- node[left,font=\scriptsize,text=textoclaro]{Sí}  (inc);
        \draw[->, acentoclaro]   (dv)   -- node[right,font=\scriptsize,text=textoclaro]{No} (dec1);
        \draw[->, acentoclaro]   (dv2)  -- node[left,font=\scriptsize,text=textoclaro]{Sí}  (dec2);
        \draw[->, acentoclaro]   (dv2)  -- node[right,font=\scriptsize,text=textoclaro]{No} (inc2);
    \end{tikzpicture}
}

\block{6. Conclusiones}{
    \begin{itemize}
        \item El convertidor Boost fue dimensionado correctamente para 8540 W.
        \item El algoritmo P\&O logra seguimiento del MPP con eficiencia del 94.3\%.
        \item El rizado de tensión se mantuvo por debajo del 2\% en régimen permanente.
        \item La simulación valida el diseño teórico propuesto \cite{mppt2020}.
    \end{itemize}
}

\block{Referencias}{
    \small
    \bibliographystyle{apacite}
    \bibliography{referencias}
}

\end{columns}
\end{document}